\section{Results}
Figure~\ref{fig:split} shows the temporal split used in the experiment.
Figure~\ref{fig:trends} shows the fitted trends before the test period.
Figure~\ref{fig:validation} shows validation curves for the selector-based
methods. Figure~\ref{fig:forecasts} displays test forecasts, while
Figures~\ref{fig:rmse} and~\ref{fig:smape} summarize test errors. The realized
roughness diagnostic is shown in Figure~\ref{fig:roughness}.

\begin{figure}[htbp]
  \centering
  \includegraphics[width=\linewidth]{\paperroot figures/sp500_train_val_test_split.pdf}
  \caption{Temporal train, validation, and test split for S\&P 500 log prices.}
  \label{fig:split}
\end{figure}

Figure~\ref{fig:split} documents the chronological evaluation design. The
training segment is used to fit candidate models, the validation segment is
used only for hyperparameter selection, and the test segment is held out until
final evaluation. This ordering is important because randomly mixing financial
time points would leak future information into the model-selection step.

\begin{figure}[htbp]
  \centering
  \includegraphics[width=\linewidth]{\paperroot figures/sp500_smoothed_trends.pdf}
  \caption{In-sample trend estimates fitted before the test period.}
  \label{fig:trends}
\end{figure}

Figure~\ref{fig:trends} compares the fitted trends on the train-plus-validation
sample before test forecasting begins. Curves that stay close to the observed
series have lower in-sample smoothing bias but may carry more short-run
variation, while flatter curves impose stronger smoothness. The figure is
therefore a visual complement to the parameter and realized-roughness tables:
it shows what each smoothing choice does to the observed log-price path.

\begin{figure}[htbp]
  \centering
  \includegraphics[width=\linewidth]{\paperroot figures/sp500_validation_curves.pdf}
  \caption{Validation curves over the Guerrero smoothness index for selector-based penalized methods. The vertical scale is clipped to keep the lower-order curves readable when the order-three curve is much larger.}
  \label{fig:validation}
\end{figure}

Figure~\ref{fig:validation} shows how validation error changes as the Guerrero
smoothness index varies for each candidate difference order. The selected point
is the minimum of the displayed validation criterion for each selector. In this
run, the order-three curves are much larger than the order-one and order-two
curves, so the plotted y-axis is limited to reveal the smaller-scale structure
that determines the useful comparison among the better-performing candidates.

\begin{figure}[htbp]
  \centering
  \includegraphics[width=\linewidth]{\paperroot figures/sp500_forecasts_test.pdf}
  \caption{Test-set forecasts from the compared methods. Pale line segments show each method's fitted train-validation trend with alpha 0.2; bright segments show the held-out test forecasts.}
  \label{fig:forecasts}
\end{figure}

Figure~\ref{fig:forecasts} overlays each method's test-period forecast on the
held-out observed log prices. The pale pre-test segments show the fitted trend
state from which each method extrapolates, which explains why the bright
forecast segments do not necessarily begin at the last raw observed value. This
figure should be read together with the test-error table: a forecast that
visually tracks the level and slope of the test period tends to have lower RMSE
and SMAPE, while forecasts that remain too flat or extrapolate the wrong slope
accumulate larger errors.

\begin{figure}[htbp]
  \centering
  \begin{subfigure}{0.48\linewidth}
    \includegraphics[width=\linewidth]{\paperroot figures/sp500_test_rmse_barplot.pdf}
    \caption{RMSE}
    \label{fig:rmse}
  \end{subfigure}
  \begin{subfigure}{0.48\linewidth}
    \includegraphics[width=\linewidth]{\paperroot figures/sp500_test_smape_barplot.pdf}
    \caption{SMAPE}
    \label{fig:smape}
  \end{subfigure}
  \caption{Test error summaries.}
\end{figure}

Figures~\ref{fig:rmse} and~\ref{fig:smape} rank methods by two test-set error
summaries. RMSE emphasizes larger forecast mistakes because errors are squared
before averaging, while SMAPE reports a scale-free percentage-style error. A
method that performs well in both panels is more robustly accurate on the
chosen test window than a method that improves only one metric.

\begin{figure}[htbp]
  \centering
  \includegraphics[width=0.8\linewidth]{\paperroot figures/sp500_roughness_vs_error.pdf}
  \caption{Realized roughness diagnostic against test RMSE.}
  \label{fig:roughness}
\end{figure}

Figure~\ref{fig:roughness} separates forecast performance from realized
roughness. The horizontal axis is the diagnostic
\(R_d(\hat{\tau})=\|D^d\hat{\tau}\|_2^2\), computed after fitting each trend;
it is not Guerrero smoothness. Points far to the left are smoother in realized
finite-difference terms, while lower points have better test RMSE. The scatter
therefore shows whether extra smoothness coincides with better forecasting in
this empirical split.

\subsection{Tables}
Table~\ref{tab:test_metrics} reports the final held-out test results. The
columns MAE, RMSE, and SMAPE summarize forecast accuracy on observations that
were not used for hyperparameter selection. The parameter columns should be
read method by method: Guerrero smoothness and \(\lambda\) are meaningful only
for finite-difference penalized smoothers, while window, alpha, and degree
describe the moving-average, exponential-smoothing, and polynomial baselines.

\begin{table}[htbp]
  \centering
  \caption{Held-out test accuracy and method-specific settings.}
  \label{tab:test_metrics}
  \begin{center}
\resizebox{\linewidth}{!}{%
\begin{tabular}{lllllllllll}
\toprule
model & n\_obs & MAE & RMSE & SMAPE & best\_order & guerrero\_smoothness & best\_lambda & window & alpha & degree \\
\midrule
TrainValidationSelector & 827 & 0.127263 & 0.148216 & 1.475928 & 1.0 & 0.05 & 0.027008 &  &  &  \\
TimeWeightedValidationSelector & 827 & 0.147868 & 0.170415 & 1.717913 & 1.0 & 0.98 & 615.756247 &  &  &  \\
HPTrend & 827 & 0.675167 & 0.774414 & 7.389961 & 2.0 & 0.944195 & 1600.0 &  &  &  \\
WhittakerTrend & 827 & 0.039826 & 0.04672 & 0.462446 & 2.0 & 0.75 & 4.874751 &  &  &  \\
MovingAverageTrend & 827 & 0.322219 & 0.367195 & 3.792746 &  &  &  & 60.0 &  &  \\
ExponentialSmoothingTrend & 827 & 0.290947 & 0.338309 & 3.416307 &  &  &  &  & 0.2 &  \\
PolynomialTrendForecaster & 827 & 0.083899 & 0.09498 & 0.973113 &  &  &  &  &  & 2.0 \\
\bottomrule
\end{tabular}%
}
\end{center}

\end{table}

Table~\ref{tab:selection_metrics} reports validation-period performance before
the final test refit. These values explain why the validation-based selectors
chose their reported settings, but they are not final generalization errors.
Comparing this table with Table~\ref{tab:test_metrics} shows whether a method
that looked good during validation continued to perform well on the held-out
test window.

\begin{table}[htbp]
  \centering
  \caption{Validation-period accuracy used for model-selection comparison.}
  \label{tab:selection_metrics}
  \begin{center}
\resizebox{\linewidth}{!}{%
\begin{tabular}{lllllllllll}
\toprule
model & n\_obs & MAE & RMSE & SMAPE & best\_order & guerrero\_smoothness & best\_lambda & window & alpha & degree \\
\midrule
TrainValidationSelector & 825 & 0.096959 & 0.118151 & 1.176689 & 1.0 & 0.05 & 0.027008 &  &  &  \\
TimeWeightedValidationSelector & 825 & 0.103816 & 0.126905 & 1.260935 & 1.0 & 0.98 & 615.756247 &  &  &  \\
HPTrend & 825 & 0.627408 & 0.742694 & 7.196558 & 2.0 & 0.944195 & 1600.0 &  &  &  \\
WhittakerTrend & 825 & 0.939478 & 1.099213 & 10.529626 & 2.0 & 0.75 & 4.874751 &  &  &  \\
MovingAverageTrend & 825 & 0.25886 & 0.290361 & 3.17557 &  &  &  & 60.0 &  &  \\
ExponentialSmoothingTrend & 825 & 0.237732 & 0.269659 & 2.911698 &  &  &  &  & 0.2 &  \\
PolynomialTrendForecaster & 825 & 0.133972 & 0.161477 & 1.629947 &  &  &  &  &  & 2.0 \\
\bottomrule
\end{tabular}%
}
\end{center}

\end{table}

Table~\ref{tab:model_parameters} collects the selected or fixed model settings
in one place. The selector rows show the chosen difference order, Guerrero
smoothness index, and corresponding \(\lambda\). HPTrend and WhittakerTrend
show their fixed finite-difference penalty settings and implied Guerrero
indices. The non-Guerrero benchmarks deliberately have blank Guerrero
smoothness entries, because their tuning parameters are window length,
exponential alpha, or polynomial degree rather than \(S_d(\lambda;N)\).

\begin{table}[htbp]
  \centering
  \caption{Selected or fixed model parameters.}
  \label{tab:model_parameters}
  \begin{center}
\resizebox{\linewidth}{!}{%
\begin{tabular}{llllllll}
\toprule
model & kind & best\_order & guerrero\_smoothness & best\_lambda & degree & window & alpha \\
\midrule
TrainValidationSelector & selector & 1.0 & 0.05 & 0.027008 &  &  &  \\
TimeWeightedValidationSelector & selector & 1.0 & 0.98 & 615.756247 &  &  &  \\
HPTrend & estimator & 2.0 & 0.944195 & 1600.0 &  &  &  \\
WhittakerTrend & estimator & 2.0 & 0.75 & 4.874751 &  &  &  \\
MovingAverageTrend & estimator &  &  &  &  & 60.0 &  \\
ExponentialSmoothingTrend & estimator &  &  &  &  &  & 0.2 \\
PolynomialTrendForecaster & polynomial\_forecaster &  &  &  & 2.0 &  &  \\
\bottomrule
\end{tabular}%
}
\end{center}

\end{table}

Table~\ref{tab:roughness_metrics} reports realized roughness for every fitted
trend using the stated roughness order. Unlike Guerrero smoothness, this
diagnostic is computed from the fitted trend itself and can therefore be
reported for all methods. The accompanying test RMSE column makes it possible
to compare the smoothness of the fitted curve with its forecasting accuracy.

\begin{table}[htbp]
  \centering
  \caption{Realized roughness diagnostics and test RMSE.}
  \label{tab:roughness_metrics}
  \begin{center}
\resizebox{\linewidth}{!}{%
\begin{tabular}{llll}
\toprule
model & roughness\_order & realized\_roughness & test\_RMSE \\
\midrule
TrainValidationSelector & 1 & 0.372955 & 0.148216 \\
TimeWeightedValidationSelector & 1 & 0.002511 & 0.170415 \\
HPTrend & 2 & 0.000126 & 0.774414 \\
WhittakerTrend & 2 & 0.010925 & 0.04672 \\
MovingAverageTrend & 2 & 0.000254 & 0.367195 \\
ExponentialSmoothingTrend & 2 & 0.018961 & 0.338309 \\
PolynomialTrendForecaster & 2 & 0.0 & 0.09498 \\
\bottomrule
\end{tabular}%
}
\end{center}

\end{table}
